\documentclass[10pt,a4paper]{amsart}
\usepackage[margin=19mm]{geometry}
\usepackage{amsmath,amssymb,mathtools}
\usepackage[scr=rsfs,cal=euler]{mathalfa}
\usepackage{xfrac}
\usepackage[unicode,hidelinks,bookmarks=false]{hyperref}
\newcommand{\ie}{i.e.\ }
\newcommand{\cf}{cf.\ }
\newcommand{\resp}{respectively\ }
\newcommand{\Subsection}{\subsection}
\providecommand{\nobreakdash}{\nobreak\hbox{-}}
\setlength{\emergencystretch}{2em}
\multlinegap0pt
\usepackage{bm}
\usepackage{bbm}
\usepackage{dsfont}
\usepackage{xspace}
\usepackage{cancel}
\usepackage{mathtools}
\usepackage{amsthm}
\makeatletter
\@ifundefined{theorem}{\newtheorem{theorem}{Theorem}}{}
\@ifundefined{lemma}{\newtheorem{lemma}[theorem]{Lemma}}{}
\makeatother
\newcommand{\dd}{\,\mathrm d}
\title[Pure-swirl loss of boundedness under $L^1\_tL^2\_x$ forcing: e]{Pure-swirl loss of boundedness under $L^1\_tL^2\_x$ forcing: exact mixed-norm ranges}
\author{Hugo Beir\~ao da Veiga, Jiaqi Yang}
\date{Source: arXiv:2608.11553v1 (CC BY 4.0)}
\begin{document}
\maketitle

\noindent\emph{Source and licence.} Pure-swirl loss of boundedness under $L^1_tL^2_x$ forcing: exact mixed-norm ranges. Version arXiv:2608.11553v1 (CC BY 4.0), distributed under the Creative Commons Attribution 4.0 International licence, as recorded in the arXiv licence metadata for this exact version and shown on the abstract page https://arxiv.org/abs/2608.11553v1. Full rights notice: licenses/arxiv-2608.11553-CC-BY-4.0.txt.\par\smallskip

\noindent\textit{Editorial scope.} Counted unit: the Lemma whose source label is \texttt{lem:velocity-rates} in the article (source0.tex lines 579--600, source label \texttt{lem:velocity-rates}), delivered together with the article's own complete original proof of it (source0.tex lines 602--664, one \texttt{proof} environment of 63 lines). No mathematics is written, repaired or re-derived: statement and proof are the article's own text line for line. Required context retained uncounted: eq:phi-pointwise (lines 451--456); eq:outer-phi (lines 464--470). Every definition, display and label of those lines is retained unchanged. Omitted from this package and not redistributed: 1--450, 457--463, 471--578, 601--601, 665--1016. Nothing inside a retained excerpt is omitted or altered: every retained body is delivered exactly as the article's lines are written, so no omission receipt of any kind accompanies this package. Citations: the retained excerpts carry no citation command at all, so no bibliography entry of the article is transcribed and no reference is redistributed. Reference adaptation: none was needed and none was made; every reference inside the delivered unit resolves inside it, so no unresolved reference or citation is produced. Printed slips kept literal: no printed word, symbol, hypothesis or number of the counted statement or of its proof was altered. Layout and class: the article's own class is \texttt{amsart} with packages including amsmath, amssymb, amsthm, mathtools, mathrsfs, microtype, hyperref, cleveref; the wrapper is the lane's pinned \texttt{amsart} editorial wrapper at 10pt on A4, which restates the theorem-like environments the retained text needs and adds the article's own macro definitions that the retained text needs (\texttt{\textbackslash dd}), with extra wrapper packages (bm, bbm, dsfont, xspace, cancel, mathtools). The delivered numbering is the unit's own. Assets: no retained span contains a figure environment, a graphics-inclusion command, a PDF-page inclusion, a table or a TikZ picture; no image file is copied into this package, the package declares no asset dependency and no image file is redistributed with it. Licence: version arXiv:2608.11553v1 of this article is distributed under the Creative Commons Attribution 4.0 International licence, as recorded in the arXiv licence metadata for this exact version and shown on the abstract page https://arxiv.org/abs/2608.11553v1. A full rights notice is delivered at licenses/arxiv-2608.11553-CC-BY-4.0.txt. Characters: every retained excerpt is pure ASCII, so the delivered engine, XeLaTeX with its default Latin Modern text font, reports no missing character.
\begin{equation}
\label{eq:phi-pointwise}
    |\phi(r,t)|\leq C\frac r{r^2+T-t},
    \qquad
    |\partial_r\phi(r,t)|\leq\frac C{r^2+T-t}.
\end{equation}

\begin{equation}
\label{eq:outer-phi}
    \phi(r,t)=-\frac\beta r,
    \quad
    \partial_r\phi(r,t)=\frac\beta{r^2}
    \quad\text{for }r\geq bR(t).
\end{equation}

\begin{lemma}[Velocity and energy rates]
\label{lem:velocity-rates}
There exist $c,C>0$ and $t_0<T$ such that, for $t_0<t<T$,
\begin{equation}
\label{eq:velocity-rate}
    c(T-t)^{-\frac{1-\alpha}{2}}
    \leq\|v(t)\|_{L^\infty(D)}
    \leq C(T-t)^{-\frac{1-\alpha}{2}}
\end{equation}
and
\begin{equation}
\label{eq:gradient-rate}
    c(T-t)^{\alpha-1}
    \leq\|\nabla v(t)\|_{L^2(D)}^2
    \leq C(T-t)^{\alpha-1}.
\end{equation}
Moreover,
\begin{equation}
\label{eq:L2-uniform}
    \sup_{0\leq t<T}\|v(t)\|_{L^2(D)}<\infty.
\end{equation}
\end{lemma}

\begin{proof}
From \eqref{eq:phi-pointwise},
\[
    |W(r,t)|\leq C\frac{r^{\alpha+1}}{r^2+s}+Cr,
    \qquad s=T-t.
\]
Writing $r=\sqrt{s}\,y$ shows that
\[
    \sup_{r\in[0,1]}\frac{r^{\alpha+1}}{r^2+s}
    \leq Cs^{(\alpha-1)/2},
\]
because $y^{\alpha+1}/(1+y^2)$ is bounded for $0<\alpha<1$.  This proves the
upper bound in \eqref{eq:velocity-rate}.  For any fixed $\rho_0>b$, put
$r(t)=\rho_0R(t)$.  For $t$ close to $T$, $r(t)<1$, and
\eqref{eq:outer-phi} yields
\[
    W(r(t),t)
    =-\beta\rho_0^{\alpha-1}R(t)^{\alpha-1}
      +\beta\rho_0R(t).
\]
The first term dominates, proving the lower bound.

For a pure swirl independent of $x_3$,
\begin{equation}
\label{eq:grad-swirl}
    |\nabla v|^2=|\partial_rW|^2+\frac{|W|^2}{r^2}.
\end{equation}
The pointwise bounds on $\phi$ and $\partial_r\phi$ give
\[
    |\partial_rW|+\frac{|W|}{r}
    \leq C\frac{r^\alpha}{r^2+s}+C.
\]
Consequently,
\begin{align}
\|\nabla v(t)\|_2^2\leq C\int_0^1\frac{r^{2\alpha+1}}{(r^2+s)^2}\,\dd r+C\leq Cs^{\alpha-1}
      \int_0^\infty\frac{y^{2\alpha+1}}{(1+y^2)^2}\,\dd y+C
     \leq Cs^{\alpha-1}.
\end{align}
For the reverse inequality, in the outer region one has
\[
    W(r,t)=\beta(r-r^{\alpha-1}).
\]
Choosing a fixed small $r_0>0$, the term $|W|/r$ is bounded below by
$cr^{\alpha-2}$ on $bR(t)\leq r\leq r_0$ for all $t$ close enough to $T$.
Thus
\[
    \|\nabla v(t)\|_2^2
    \geq c\int_{bR(t)}^{r_0}r^{2\alpha-3}\,\dd r
    \geq cs^{\alpha-1},
\]
which proves \eqref{eq:gradient-rate}.

Finally,
\begin{align}
    \|v(t)\|_2^2
    &\leq C\int_0^1\frac{r^{2\alpha+3}}{(r^2+s)^2}\,\dd r+C\\
    &=Cs^\alpha\int_0^{1/\sqrt{s}}
       \frac{y^{2\alpha+3}}{(1+y^2)^2}\,\dd y+C
      \leq C,
\end{align}
where the last estimate follows by splitting the $y$-integral at $1$.
This proves \eqref{eq:L2-uniform}.
\end{proof}

\end{document}
